\section{Mở ma trận ra xem}
\label{sec:ch06-mo-ma-tran}

\index{ma trận đồng chỉnh!vì sao đọc được}%
Mọi mục trước đo mô hình từ bên ngoài: đưa câu vào, đọc câu ra, đếm câu đúng.
Cơ chế này cho phép làm một việc mà chương~\ref{ch:encoder-decoder} không làm
được. Các \(\alpha_{ij}\) là số thật, nằm sẵn trong mô hình, và xếp chúng thành
bảng thì được một \tn{ma trận đồng chỉnh}{alignment matrix}: hàng \(i\) là câu
trả lời cho \enquote{lúc viết từ đích thứ \(i\), mô hình đọc những chỗ nào của
câu nguồn}.

Bài dành mục 5.2.1 cho việc ấy và in bốn ma trận Anh-Pháp. Nhận xét của họ là
\tn{đồng chỉnh}{alignment} phần lớn đi theo đường chéo, với vài chỗ chệch ra ở
tính từ và danh
từ, vì hai thứ tiếng xếp chúng khác nhau.

Corpus của sách này dựng sao cho chỗ chệch ấy không phải ngoại lệ mà là toàn bộ
bài toán. Tiếng Anh viết \code{black cat}, tiếng Việt viết \code{mèo đen}. Mọi
tính từ trong corpus đều phải bắt chéo qua danh từ của nó.

\subsection*{Một câu, tám hàng}

\begin{measured}{ma trận đồng chỉnh của câu \code{the black cat wants a new lamp}}
\begin{minted}{text}
target        1     2     3     4     5     6     7
con        0.07  0.15  0.39  0.38  0.01  0.00  0.00
mèo        0.00  0.04  0.95  0.01  0.00  0.00  0.00
đen        0.00  0.82  0.15  0.03  0.00  0.00  0.00
muốn       0.00  0.00  0.00  0.99  0.01  0.00  0.00
một        0.00  0.00  0.00  0.05  0.95  0.00  0.00
cái        0.00  0.00  0.00  0.00  0.00  0.05  0.95
đèn        0.00  0.00  0.00  0.00  0.00  0.18  0.82
mới        0.00  0.00  0.00  0.00  0.00  0.98  0.02
\end{minted}

Cột 1 tới 7 là \code{the black cat wants a new lamp}. \(d_h = 32\), 14 epoch,
câu nguồn \emph{không} đảo, khớp đúng cả câu 0.9967 trên tập test.

Đo bằng \code{experiments/ch06\_alignment.py} tại tag \code{ch06}.
\end{measured}

Bảng này là mô hình tự khai, nên đọc chậm được.

Hàng \code{mèo} dồn 0.95 vào cột 3, tức \code{cat}. Hàng ngay sau,
\code{đen}, dồn 0.82 vào cột 2, tức \code{black}. Chỉ số cột \emph{giảm} khi đi
từ hàng này xuống hàng sau. Một đồng chỉnh đường chéo không làm được chuyện đó,
vì đường chéo thì chỉ số cột chỉ có tăng.

Cặp thứ hai lặp lại y hệt ở cuối câu: \code{đèn} nhìn cột 7 (\code{lamp}) với
0.82, rồi \code{mới} nhìn cột 6 (\code{new}) với 0.98.

Ở giữa, \code{muốn} nhìn \code{wants} với 0.99 và \code{một} nhìn \code{a} với
0.95. Hai hàng ấy đi theo đường chéo, đúng như phải thế.

Hình~\ref{fig:ch06-ma-tran} vẽ lại cùng bảng ấy theo quy ước xám của bài.

\begin{figure}[htbp]
  \centering
  % Ma trận đồng chỉnh đo được, tag ch06, experiments/ch06_alignment.py.
% Hàng là từ tiếng Việt mô hình sinh ra, cột là từ tiếng Anh của câu nguồn.
% Mức xám là trọng số alpha_ij: trắng là 0, đen là 1. Vẽ theo hình 3 của
% Bahdanau 2015, cùng quy ước xám.
% Mono-safe theo cấu tạo: thông tin nằm hoàn toàn ở mức xám, không ở màu.
\begin{tikzpicture}[
  o/.style     = {draw = black!30, line width = 0.2pt, minimum size = 6.2mm,
                  inner sep = 0pt},
  nhan/.style  = {font = \scriptsize},
  cot/.style   = {font = \scriptsize, rotate = 55, anchor = west},
]

  % Cột: câu nguồn tiếng Anh.
  \foreach \j/\lab in {0/the, 1/black, 2/cat, 3/wants, 4/a, 5/new, 6/lamp} {
    \node[cot] at (\j*0.62, 0.42) {\lab};
  }

  % Mỗi hàng: nhãn tiếng Việt rồi bảy trọng số, đọc thẳng từ output của script.
  \foreach \r/\lab/\va/\vb/\vc/\vd/\ve/\vf/\vg in {%
    0/con/7/15/39/38/1/0/0,
    1/mèo/0/4/95/1/0/0/0,
    2/đen/0/82/15/3/0/0/0,
    3/muốn/0/0/0/99/1/0/0,
    4/một/0/0/0/5/95/0/0,
    5/cái/0/0/0/0/0/5/95,
    6/đèn/0/0/0/0/0/18/82,
    7/mới/0/0/0/0/0/98/2%
  } {
    \node[nhan, anchor = east] at (-0.42, -\r*0.62) {\lab};
    \foreach \j/\v in {0/\va, 1/\vb, 2/\vc, 3/\vd, 4/\ve, 5/\vf, 6/\vg} {
      \node[o, fill = black!\v] at (\j*0.62, -\r*0.62) {};
    }
  }

  % Hai chỗ bắt chéo, khoanh bằng nét đứt chứ không bằng màu.
  % Ô (cột j, hàng r) tâm ở (0.62j, -0.62r) và cạnh 0.62, nên một khung ôm
  % các cột j1..j2 và các hàng r1..r2 chạy từ (0.62*j1 - 0.31, -0.62*r1 + 0.31)
  % tới (0.62*j2 + 0.31, -0.62*r2 - 0.31). Khung thứ nhất ôm mèo-đen với
  % black-cat, khung thứ hai ôm đèn-mới với new-lamp.
  \draw[black!70, densely dashed, line width = 0.6pt, rounded corners = 1pt]
    (0.31, -0.31) rectangle (1.55, -1.55);
  \draw[black!70, densely dashed, line width = 0.6pt, rounded corners = 1pt]
    (2.79, -3.41) rectangle (4.03, -4.65);

\end{tikzpicture}

  \caption{Cùng ma trận, vẽ theo quy ước hình 3 của bài: trắng là 0, đen là 1.
    Hai khung nét đứt khoanh hai chỗ bắt chéo. Đọc dọc theo phần đậm thì thấy
    nó đi xuống theo đường chéo rồi giật ngược lại một cột, đúng hai lần, đúng
    hai chỗ mà tiếng Việt đặt tính từ sau danh từ. Hàng đầu, \code{con}, là
    hàng duy nhất không có ô nào thật đậm.}
  \label{fig:ch06-ma-tran}
\end{figure}

\subsection*{Hàng đầu tiên là hàng thú vị nhất}

\index{loại từ!attention tản trên nhiều vị trí}%
\code{con} là \tn{loại từ}{classifier}, và nó không ứng với từ tiếng Anh nào.
Trong \code{the black cat} không có token nào để nó nhìn.

Nên nó tản: 0.39 vào \code{cat}, 0.38 vào \code{wants}, 0.15 vào \code{black},
0.07 vào \code{the}. Không có ô nào vượt 0.4.

Đây đúng là tình huống bài mô tả ở mục 5.2.1, và họ nói về nó bằng ví dụ tiếng
Pháp. \tn{Đồng chỉnh cứng}{hard alignment} buộc mỗi từ đích phải chỉ vào một từ
nguồn, nên nó phải bịa ra một chỗ để chỉ, hoặc phải chỉ vào một ký hiệu rỗng mà
bài gọi là \code{[NULL]}. \tn{Đồng chỉnh mềm}{soft alignment} không cần thế: nó
xử lý được cụm nguồn và cụm đích khác độ dài mà không phải ánh xạ từ nào tới
hoặc từ hư không. Nguyên văn: \enquote{without requiring a counter-intuitive way
of mapping some words to or from nowhere}.

Loại từ tiếng Việt là ví dụ sạch nhất cho câu ấy mà tôi biết. Nó là một từ bắt
buộc phải có ở \tn{câu đích}{target sentence}, tương ứng với không gì cả ở câu
nguồn, và mô hình xử
lý nó bằng cách nhìn mờ vào cả cụm.

\index{loại từ!hai loại từ hành xử khác nhau}%
Rồi so \code{con} với \code{cái}, cũng là loại từ, ở hàng thứ sáu. \code{cái}
dồn 0.95 vào \code{lamp}. Nó không tản chút nào.

Hai từ cùng loại và hành xử ngược nhau. Cách đọc tôi thấy hợp lý nhất là loại
từ tiếng Việt phụ thuộc vào danh từ đi sau nó, \code{con} cho vật sống và
\code{cái} cho đồ vật, nên chỗ đáng nhìn để chọn loại từ chính là danh từ. Vậy
\code{cái} nhìn \code{lamp} là hành vi đúng, còn \code{con} lẽ ra cũng nên dồn
vào \code{cat} như thế.

Tôi không có phép đo nào tách được hai khả năng: hoặc hàng \code{con} là chỗ mô
hình còn chưa học xong, hoặc vị trí đầu câu có gì đó khác. Nó là hàng đầu tiên,
mà bước đầu tiên là bước duy nhất chấm điểm bằng \(\vect{s}_0\) dựng từ
encoder chứ không từ một bước decoder nào trước đó. Đó là một
\tn{giả thuyết}{hypothesis}, tôi chưa kiểm, và bài tập cuối chương giao lại
việc kiểm.

\subsection*{Một câu không phải bằng chứng}

\index{ma trận đồng chỉnh!đo tỷ lệ bắt chéo}%
Mọi thứ trên đây đọc từ một ma trận. Bài cũng đọc từ bốn ma trận, và cách đọc
ấy là cách đọc bằng mắt: nó thuyết phục đúng chừng mực của một ví dụ được chọn.

Corpus này cho phép làm chặt hơn, vì chỗ bắt chéo ở đây định nghĩa được. Duyệt
toàn bộ 300 câu test. Trong output của \emph{chính mô hình}, tìm mọi chỗ một
danh từ tiếng Việt đi liền trước một tính từ. Với mỗi cặp ấy, lấy vị trí nguồn
mà mỗi từ dồn trọng số nhiều nhất, rồi hỏi vị trí của tính từ có nằm bên
\emph{trái} vị trí của danh từ không.

Một đồng chỉnh đường chéo cho 0 phần trăm ở phép đo này, theo cấu tạo.

\begin{measured}{tỷ lệ bắt chéo trên toàn bộ tập test}
\begin{minted}{text}
noun-adjective pairs in the model's own output: 625
of those, adjective attends left of the noun: 609
crossing rate: 0.9744
\end{minted}

Cùng mô hình với ma trận ở trên.

Đo bằng \code{experiments/ch06\_alignment.py} tại tag \code{ch06}.
\end{measured}

625 cặp, 609 bắt chéo, tỷ lệ 0.9744. Phép đo đi theo output của mô hình chứ
không theo câu tham chiếu, nên một câu mô hình dịch sai vẫn đóng góp nếu nó
sinh ra một cặp danh từ và tính từ. Cái được đo là mô hình \emph{nhìn} vào đâu,
không phải nó có đúng không.

\subsection*{Một ma trận đẹp không chứng minh bản dịch tốt hơn}

\index{ma trận đồng chỉnh!không phải bằng chứng chất lượng}%
Chỗ này cần nói ra, vì cả mục vừa rồi mời gọi đúng cái kết luận sai ấy.

Luong và cộng sự đo thẳng chuyện đó một năm sau. Họ có dữ liệu đồng chỉnh chuẩn
do người làm cho 508 câu Anh-Đức, nên họ chấm được alignment error rate của các
mô hình. Kết quả: ensemble của họ đạt 0.34, bằng đúng mô hình local-m đơn lẻ,
trong khi ensemble dịch tốt hơn hẳn. Họ rút ra nhận xét thẳng thắn rằng
\tn{điểm đồng chỉnh}{alignment score} và điểm dịch không tương quan tốt với
nhau (\enquote{are not well correlated}), và dẫn Fraser và
Marcu 2007 cho nhận xét ấy~\autocite{luong2015effective}.

Nên bảng 0.9744 ở trên nói rằng mô hình đã học được một quy tắc thật về hai thứ
tiếng, và \emph{không} nói rằng nó dịch tốt hơn vì thế. Cái nói chuyện thứ hai
là bảng ở mục~\ref{sec:ch06-be-rong}, và hai bảng ấy độc lập với nhau.

Cách tôi dùng ma trận đồng chỉnh, sau khi đọc nhận xét của Luong, là dùng nó
như một công cụ chẩn đoán chứ không như một thước đo chất lượng. Nó trả lời
được câu \enquote{mô hình đang nhìn vào đâu}, và câu đó đắt giá lúc gỡ lỗi.
Chương~\ref{ch:encoder-decoder} không có cách nào hỏi câu ấy.
